\documentclass[11pt,a4paper]{article}

\usepackage[utf8]{inputenc}
\usepackage[T1]{fontenc}
\usepackage{lmodern}
\usepackage{amsmath,amssymb,amsthm,mathtools}
\usepackage{booktabs}
\usepackage{enumitem}
\usepackage{microtype}
\usepackage{xcolor}
\usepackage{geometry}
\usepackage{hyperref}
\geometry{margin=1in}
\hypersetup{
  colorlinks=true,
  linkcolor=blue!55!black,
  citecolor=blue!55!black,
  urlcolor=blue!55!black
}

\newtheorem{theorem}{Theorem}[section]
\newtheorem{corollary}[theorem]{Corollary}
\newtheorem{proposition}[theorem]{Proposition}
\theoremstyle{definition}
\newtheorem{definition}[theorem]{Definition}
\newtheorem{assumption}[theorem]{Assumption}
\theoremstyle{remark}
\newtheorem{remark}[theorem]{Remark}

\newcommand{\dd}{\mathrm d}
\newcommand{\cP}{\mathcal P}
\newcommand{\cQ}{\mathcal Q}
\newcommand{\cE}{\mathcal E}
\newcommand{\cH}{\mathcal H}
\newcommand{\cO}{\mathcal O}
\newcommand{\TT}{\mathrm{TT}}
\newcommand{\eff}{\mathrm{eff}}

\title{\textbf{Controlled Coherent Reduction to\\
Four-Dimensional Einstein Gravity}}
\author{Peter Nero}
\date{Third edition\\July 2026}

\begin{document}
\maketitle

\begin{abstract}
This paper identifies the precise sense in which a Modal Triplet Theory (MTT)
realization can reproduce classical four-dimensional General Relativity.  The
starting point is a ten-dimensional bundle over a Lorentzian four-manifold
with compact six-dimensional fiber \(X_6\), together with a local covariant
action whose gravitational sector contains the higher-dimensional
Einstein--Hilbert term.  That metric and action are realization inputs; they
are not derived here from projection alone.  We formulate an exact
consistent-truncation condition and an approximate gapped version that bounds
discarded modes and the induced effective-action error.  When all unwanted
metric zero modes are absent, stabilized, or retained explicitly, the
two-derivative metric sector reduces to Einstein--Hilbert form with
\(\kappa_4^2=\kappa_{10}^2/V_W\), where \(V_W\) is the warped internal volume.
Diffeomorphism invariance then gives the Einstein equation and covariant
stress-energy conservation; universal minimal coupling supplies the usual
geometric-optics and test-body limits.  The current MTT calculation also
provides an exact finite-branch transverse-traceless helicity-two support
certificate.  It does not yet determine the physical Newton/Planck
normalization, the complete stress-energy response, the cosmological
constant, or a projection-only source for the Lorentzian action.  The result
is therefore a controlled GR-limit theorem and a sharply delimited research
bridge, not a first-principles derivation of gravity from MTT axioms alone.
\end{abstract}

\section*{Revision note for this edition}
\addcontentsline{toc}{section}{Revision note for this edition}
\begin{description}[leftmargin=3.5cm,style=nextline]
\item[Supersedes.] \emph{Modal Triplet Theory: From MTT to General
Relativity}, version~2.
\item[Reason.] The previous edition called a standard dimensional reduction a
first-principles derivation even though it inserted the ten-dimensional
Einstein--Hilbert action at the outset.  It also wrote
\(Y_4\times B_1\times B_2\times B_3\) with three three-dimensional factors
and called the result ten-dimensional, omitted the nonlinear
consistent-truncation test, and suppressed possible scalar and vector zero
modes.
\item[Resolution.] This edition uses \(M_{10}\to Y_4\) with
six-dimensional fiber \(X_6\), separates imported causal/action data from
derived reduction statements, supplies exact and approximate truncation
criteria, treats moduli and Kaluza--Klein vectors explicitly, and restricts
the uniqueness claim to the hypotheses of the four-dimensional
two-derivative metric sector.
\item[Retained result.] Given a higher-dimensional Einstein--Hilbert sector
and a valid coherent truncation, the internal pushforward does yield a
four-dimensional Einstein--Hilbert term and the standard Einstein equations
with an effective coupling fixed by the warped volume.
\item[Remaining boundary.] MTT has not yet selected the physical Lorentzian
principal symbol and action from projection alone, nor derived the SI value
of Newton's constant, the full matter stress response, or the observed
cosmological constant from the same source.  The exact finite-branch
helicity-two support result closes an internal support gate only.
\end{description}

\tableofcontents

\section{Overview: Question, Answer, and Logical Status}

\subsection{The question actually answered}

The paper asks a conditional but important question:
\begin{quote}
Given a ten-dimensional Lorentzian metric theory on a selected MTT carrier,
when does its coherent low-energy sector reduce to four-dimensional Einstein
gravity, and what is still needed before that reduction becomes an MTT source
theorem?
\end{quote}
There are four different operations in that sentence.  Keeping them separate
prevents a valid calculation from being assigned a stronger meaning than it
has:
\[
\begin{gathered}
\boxed{\text{source selects geometry and action}}
\longrightarrow
\boxed{\text{retain a coherent mode sector}}
\\[1mm]
\Downarrow
\\[-1mm]
\boxed{\text{compare the 4D observables with GR}}
\longleftarrow
\boxed{\text{integrate out discarded modes}} .
\end{gathered}
\]
This paper proves implications in the last three boxes after the first has
been supplied.  It does not prove the first box.

\subsection{Current MTT status}

\begin{center}
\begin{tabular}{p{0.24\linewidth}p{0.66\linewidth}}
\toprule
Status & Content \\
\midrule
Established internally &
The selected finite exact branch has a two-dimensional helicity-two
transverse-traceless support carrier and normalized internal eigenvalue
\(\lambda_{\mathrm{GR},\TT}=15\). \\
Conditional bridge &
A supplied ten-dimensional Lorentzian action reduces to a four-dimensional
metric action under the consistency, gap, stabilization, and locality
conditions stated below. \\
Not selected here &
The Lorentzian metric or principal symbol, the Einstein--Hilbert source
coefficient, the selected compactification solution, and the boundary and
quantization data. \\
Still open physically &
Newton/Planck normalization, complete stress-energy response, the observed
cosmological constant, and a same-source derivation connecting the internal
TT carrier to the normalized spacetime graviton. \\
\bottomrule
\end{tabular}
\end{center}

The distinction matters.  A helicity-two carrier is necessary for a graviton
sector, but it is not by itself a kinetic action, a universal coupling law, or
Einstein's nonlinear equation.

\paragraph{In plain language.}
The compact geometry can tell us which internal patterns survive and how a
supplied gravitational action is averaged.  It cannot manufacture a
Lorentzian causal law or its physical coupling merely by deleting modes.  The
paper therefore treats ``which pattern survives,'' ``which equation governs
it,'' and ``how strongly it couples'' as three separate questions.

\section{Typed Ten-Dimensional Geometry}

\subsection{Base, fiber, and the
\texorpdfstring{\(1<2<3\)}{1<2<3} filtration}

Let
\[
 \pi:M_{10}\longrightarrow Y_4
\]
be a smooth bundle whose base \(Y_4\) is a time-oriented Lorentzian
four-manifold and whose compact fiber is a six-manifold \(X_6\).  In a local
trivialization one may write \(M_{10}\simeq Y_4\times X_6\), but the global
bundle need not be a Cartesian product.

The MTT rank flag
\[
 1<2<3,\qquad 1+2+3=6,
\]
is a filtration of internal carrier data.  It is not a product of three
independent three-manifolds.  The shared circle is represented by one common
\(U(1)\) line bundle and connection reused by the relevant lanes; it is
counted once and is not identified with Lorentzian time.  The selected q79
Fu--Yau branch is the present compactification candidate.  The auxiliary
Lens--Nil model can still test local operators, but it is not interchangeable
with that global six-manifold.

\subsection{Why the Lorentzian metric is an input}

A local comparison field
\[
 Q_{\rm WW}\in\Gamma\!\left(\operatorname{Hom}(TP,TI)\right)
\]
may define a positive comparison tensor \(Q_{\rm WW}^{T}Q_{\rm WW}\).  Such a
Gram tensor is positive semidefinite and therefore cannot, without additional
structure, be the Lorentzian spacetime metric.  A physical realization must
supply a Lorentzian principal symbol, or equivalently suitable frame/coframe
and soldering data, before causal cones and hyperbolic propagation are
defined.  Projecting a positive internal norm does not change its signature
by itself.

\subsection{A useful local metric ansatz}

For the reduction calculation, take a background of the form
\begin{equation}
 \dd s_{10}^{\,2}
 =e^{2A(y)}g_{\mu\nu}(x)\dd x^\mu\dd x^\nu
  h_{mn}(y)\dd y^m\dd y^n ,
\label{eq:warpedmetric}
\end{equation}
with fixed warp factor \(A\) and fixed internal metric \(h\) in the strict
metric-only truncation.  More general off-diagonal components and
\(x\)-dependent moduli are not errors; they become four-dimensional vector
and scalar fields and must be retained or consistently stabilized.

\section{The Higher-Dimensional Action Is a Realization Input}

Consider a local covariant effective action
\begin{equation}
 S_{10}
 =\frac{1}{2\kappa_{10}^{2}}
  \int_{M_{10}}\!\sqrt{|g_{10}|}\,
  \bigl(R_{10}-2\Lambda_{10}\bigr)\,\dd^{10}X
 +S_{\rm bulk}[g_{10},\Phi]
 +S_{\rm bdy}.
\label{eq:S10}
\end{equation}
The boundary term must make the chosen variational problem well posed.
\(S_{\rm bulk}\) can contain gauge fields, spinors, scalars, flux, torsion,
and higher-derivative operators.  Equation~\eqref{eq:S10} is a declared
action ansatz in the relevant regime.  Varying it rigorously derives its field
equations; that variation does not explain why MTT selects this action rather
than another diffeomorphism-invariant action.

At energies \(E\ll\Lambda_{10}\), the omitted operators are organized as a
derivative expansion.  Terms such as
\[
 R_{10}^{2},\quad R_{AB}R^{AB},\quad
 R_{ABCD}R^{ABCD},\quad R_{10}F^2,\quad F^4
\]
are generally allowed unless a symmetry or a source theorem removes them.
Their coefficients cannot be inferred from a spectral gap alone.

\section{Coherent Modes and Consistent Truncation}

\subsection{Retained and discarded fields}

Let \(\cP\) be the selected coherent projector and
\(\cQ=1-\cP\).  Decompose every field, including metric fluctuations, as
\[
 \Phi=\Phi_{\cP}+\eta_{\cQ}.
\]
Equivalently, for normalized internal modes \(\chi_n\),
\[
 \Phi(x,y)=
 \sum_{a\in I_{\rm keep}}\phi_a(x)\chi_a(y)
 +\sum_{r\in I_{\rm disc}}\eta_r(x)\chi_r(y).
\]
A harmonic projector is useful for identifying zero modes, but harmonicity
alone does not prove that nonlinear products of retained modes fail to source
discarded ones.

\begin{definition}[Exact coherent truncation]
The retained ansatz is an exact consistent truncation when
\begin{equation}
 \left.
 \cQ\,\frac{\delta S_{10}}{\delta\Phi}
 \right|_{\eta_{\cQ}=0}=0
\label{eq:exacttrunc}
\end{equation}
for every retained configuration in the stated domain.  Then every solution
of the reduced equations uplifts to a solution of the full equations within
that ansatz.
\end{definition}

For an effective rather than exact truncation, define the discarded source
\[
 J_{\cQ}(\Phi_{\cP})
 :=
 \left.
 \cQ\,\frac{\delta S_{10}}{\delta\Phi}
 \right|_{\eta_{\cQ}=0}.
\]
This is the object that a projection calculation must bound.  The condition
\(\|J_{\cQ}\|\le\epsilon\) is the explicit form of the source-leakage gate.

\begin{theorem}[Gapped approximate coherent reduction]
\label{thm:controlled}
Fix Banach norms on retained and discarded fields in a neighborhood
\(\mathcal U\).  Suppose:
\begin{enumerate}[label=(\roman*)]
\item \(S_{10}\) is twice differentiable on \(\mathcal U\);
\item the discarded linearized operator
\[
 L_{\cQ}
 =\left.
 \cQ\,\frac{\delta^2S_{10}}{\delta\Phi^2}\,\cQ
 \right|_{\eta_{\cQ}=0}
\]
has an inverse on the gauge-fixed discarded subspace with
\(\|L_{\cQ}^{-1}\|\le m_{\rm gap}^{-2}\);
\item the nonlinear remainder in the discarded equation is Lipschitz in
\(\eta_{\cQ}\) with constant \(L_N<m_{\rm gap}^{2}\); and
\item \(\|J_{\cQ}(\Phi_{\cP})\|\le\epsilon\) uniformly.
\end{enumerate}
Then the discarded equation has a unique local solution satisfying
\begin{equation}
 \|\eta_{\cQ}\|
 \le
 \frac{\epsilon}{m_{\rm gap}^{2}-L_N}.
\label{eq:etabound}
\end{equation}
Substitution into the action gives the Schur--Feshbach correction
\begin{equation}
 S_{\eff}[\Phi_{\cP}]
 =S_{10}[\Phi_{\cP},0]
 -\frac12\langle J_{\cQ},L_{\cQ}^{-1}J_{\cQ}\rangle
 +\cO\!\left(\frac{\epsilon^3}{m_{\rm gap}^{6}}\right)
\label{eq:schur}
\end{equation}
whenever the stated expansion is uniform.  Exact consistency is the special
case \(J_{\cQ}=0\).
\end{theorem}

\begin{proof}
Write the discarded Euler--Lagrange equation as
\[
 L_{\cQ}\eta_{\cQ}+J_{\cQ}+N(\eta_{\cQ})=0 .
\]
The map
\(
 T(\eta)=-L_{\cQ}^{-1}[J_{\cQ}+N(\eta)]
\)
is a contraction on the ball of radius
\(\epsilon/(m_{\rm gap}^{2}-L_N)\).  The Banach fixed-point theorem gives the
unique solution and bound~\eqref{eq:etabound}.  Taylor expansion of the action
at \(\eta_{\cQ}=0\), followed by
\(\eta_{\cQ}=-L_{\cQ}^{-1}J_{\cQ}+\cO(\epsilon^2/m_{\rm gap}^4)\),
gives~\eqref{eq:schur}.
\end{proof}

\subsection{The zero-mode audit}

The gap in Theorem~\ref{thm:controlled} applies only on the complement of all
zero modes.  A compactification must therefore inventory them before claiming
pure GR:
\begin{itemize}
\item deformations of \(h_{mn}\) can produce scalar moduli;
\item off-diagonal metric components \(g_{\mu m}\) can produce vector fields,
especially along internal isometries or invariant one-forms;
\item form fields can produce additional scalars, vectors, and antisymmetric
tensors; and
\item gauge and diffeomorphism zero modes must be removed or treated by a
declared gauge-fixing and ghost complex.
\end{itemize}
Each unwanted physical zero mode must be absent, stabilized above the working
scale, or retained in the four-dimensional theory.  Declaring it
``incoherent'' is not a substitute for checking its equation of motion.

\section{Four-Dimensional Einstein Sector}

\subsection{Warped volume and Newton coupling}

For the metric~\eqref{eq:warpedmetric}, the coefficient of \(R_4(g)\) is
controlled by the warped volume
\begin{equation}
 V_W:=\int_{X_6}\sqrt{h}\,e^{2A(y)}\,\dd^6y .
\label{eq:VW}
\end{equation}
If the volume and warp modes are fixed, the leading four-dimensional metric
term is
\begin{equation}
 S_{\rm grav}^{(4)}
 =\frac{1}{2\kappa_4^2}
 \int_{Y_4}\sqrt{-g}\,(R_4-2\Lambda_4)\,\dd^4x,
\qquad
 \kappa_4^2=\frac{\kappa_{10}^2}{V_W}.
\label{eq:kappa4}
\end{equation}
In the unwarped direct-product case, \(V_W=\operatorname{Vol}(X_6)\).
Internal curvature, flux, vacuum energy, branes, and higher-dimensional
\(\Lambda_{10}\) all contribute to the four-dimensional potential.
Consequently \(\Lambda_4\) is not, in general, just minus one half of an
averaged internal scalar curvature.

Equation~\eqref{eq:kappa4} is a dictionary, not yet a prediction.  It predicts
Newton's constant only when both \(\kappa_{10}\) and the normalized selected
volume are derived independently.  With the reduced Planck convention,
\[
 M_{\rm Pl,red}^{2}
 =\kappa_4^{-2}
 =\frac{V_W}{\kappa_{10}^{2}}.
\]

\subsection{Why Einstein--Hilbert is the leading metric action}

Suppose the low-energy retained gravitational field is only \(g_{\mu\nu}\),
the action is local and four-dimensionally diffeomorphism invariant, and the
metric field equations contain no derivatives above second order.  Under the
standard natural-tensor assumptions, Lovelock's four-dimensional result
restricts the field equation to a linear combination of
\(G_{\mu\nu}\) and \(g_{\mu\nu}\).  Equivalently, the leading action is
Einstein--Hilbert plus a cosmological term, up to a boundary term and a
four-dimensional topological Gauss--Bonnet contribution.

This is a conditional uniqueness statement.  It does not apply unchanged
when additional light fields survive, locality is relaxed, higher derivatives
are retained, or the causal principal symbol is not the metric one.

\begin{corollary}[Controlled Einstein limit]
\label{cor:GR}
Assume the hypotheses of Theorem~\ref{thm:controlled}, stabilization or
retention of every physical zero mode, universal minimal coupling of retained
matter, and the metric-only two-derivative conditions above.  Then, throughout
the controlled regime, the leading metric equation is
\begin{equation}
 G_{\mu\nu}+\Lambda_4g_{\mu\nu}
 =\kappa_4^2T_{\mu\nu}
 +\Delta_{\mu\nu},
\label{eq:einstein}
\end{equation}
where \(\Delta_{\mu\nu}=0\) for an exact two-derivative truncation and is
bounded by the discarded-mode and higher-derivative errors in the approximate
case.
\end{corollary}

\begin{proof}
Theorem~\ref{thm:controlled} supplies a controlled local effective action.
The field-content and derivative assumptions reduce its leading metric part
to~\eqref{eq:kappa4}.  Varying with respect to \(g^{\mu\nu}\) gives
Eq.~\eqref{eq:einstein}; the remaining terms define
\(\Delta_{\mu\nu}\).
\end{proof}

\subsection{Bianchi identity and stress-energy}

For
\[
 T_{\mu\nu}
 :=-\frac{2}{\sqrt{-g}}
   \frac{\delta S_{\rm matter}^{(4)}}{\delta g^{\mu\nu}},
\]
four-dimensional diffeomorphism invariance gives the Noether identity
\[
 \nabla^\mu T_{\mu\nu}=0
\]
on the retained matter equations.  This conservation law follows from the
symmetry of the reduced action and the matter equations; it is not produced
merely by averaging over \(X_6\).  In an approximate truncation the complete
effective stress tensor, including the correction terms, is conserved.

\section{Motion, Equivalence, and Classical Tests}

If all retained matter sectors couple to the same metric, a minimally coupled
scalar in the WKB form
\(
\phi=a\,e^{iS/\varepsilon}
\)
obeys at leading order
\[
 g^{\mu\nu}\partial_\mu S\,\partial_\nu S+m^2=0.
\]
Its rays are timelike geodesics; the corresponding massless equation gives
null geodesics.  The same principal-symbol argument applies to standard
minimally coupled wave equations.  This establishes the geometric-optics
limit and the weak equivalence principle within the declared universal
coupling regime.

The statement has ordinary limits.  Spinning bodies, extended bodies,
self-gravitating objects, and fields with nonminimal curvature couplings can
obey corrected motion equations.  Projection alone does not prove universal
minimal coupling; the common matter-metric source map must supply it.

Once Corollary~\ref{cor:GR} applies with negligible
\(\Delta_{\mu\nu}\), the Newtonian, post-Newtonian, redshift, lensing, and
gravitational-wave predictions are those of GR with
\(G_N=\kappa_4^2/(8\pi)\).  This is agreement by reduction.  It is not a new
numerical prediction until the right-hand side of~\eqref{eq:kappa4} is
selected without importing \(G_N\).

\section{The Exact MTT Helicity-Two Support Result}

The current GR successor calculation establishes, on its selected finite
exact branch,
\begin{equation}
 \Pi_{\rm exact64}B^\ast P_{\TT}=B^\ast P_{\TT},
\label{eq:TTsupport}
\end{equation}
or equivalently
\begin{equation}
 \operatorname{supp}(J_{\TT})
 =|d_\ast\rangle\otimes
 \operatorname{span}\{c_2,s_2\},
\qquad
 \lambda_{\mathrm{GR},\TT}=15
\label{eq:TTlambda}
\end{equation}
in normalized internal exact-branch units.  The real
\(k=2\) and \(k=62\) character plane of \(\mathbb C[\mathbb Z_{64}]\) carries
the two same-angle helicity-two directions.  The associated certificate
checks rank two, no support leakage outside the selected carrier, and the
required finite central-shift intertwining.

\subsection{What this closes}

Equations~\eqref{eq:TTsupport}--\eqref{eq:TTlambda} close the internal
transverse-traceless support question for that exact branch.  They provide a
nontrivial compatibility test: the selected finite carrier has precisely the
two real polarizations required by a four-dimensional massless helicity-two
sector.
This result establishes carrier compatibility, not the full dynamics of the
carrier.

\subsection{What it does not close}

The support certificate does not by itself establish:
\begin{enumerate}
\item the spacetime kinetic operator and its Lorentzian principal symbol;
\item the residue and positivity of the physical graviton pole;
\item the normalization \(\kappa_4\) or the SI value of \(G_N\);
\item universal coupling to the complete conserved stress tensor;
\item the nonlinear Einstein self-coupling; or
\item the absence of extra scalar, vector, ghost, or massive spin-two modes.
\end{enumerate}
Those are distinct gates.  Conflating support with normalization was one of
the reasons the earlier GR claim was too strong.

\section{Corrections Beyond the Strict Limit}

The controlled four-dimensional effective action can contain
\begin{align}
 S_{\eff}
 =\int\sqrt{-g}\,\dd^4x\,
 \Big[
 &\frac{M_{\rm Pl,red}^{2}}{2}(R-2\Lambda_4)
 +c_1R^2+c_2R_{\mu\nu}R^{\mu\nu}
 +c_3R_{\mu\nu\rho\sigma}R^{\mu\nu\rho\sigma}
 \nonumber\\
 &+\mathcal L_{\rm light}(g,\varphi,A,\psi)+\cdots
 \Big].
\label{eq:EFT}
\end{align}
In four dimensions one curvature-squared combination is topological, leaving
two local combinations before field redefinitions and matter couplings are
considered.  The coefficients \(c_i\) depend on the supplied
higher-dimensional action, thresholds, background, and matching scheme.
They are not universally proportional to a projector gap.

Around Minkowski space, a local Lorentz-invariant higher-derivative kinetic
operator often factors schematically as
\[
 k^2\bigl(1+\alpha k^2+\cdots\bigr).
\]
The ordinary massless pole \(k^2=0\) remains luminal; the additional factor
can introduce heavy poles or signal the limit of the EFT.  It is therefore
incorrect to read the expression automatically as a frequency-dependent
speed for the massless graviton.  Background curvature, Lorentz-violating
operators, or nonlocal terms require their own principal-symbol analysis.

\section{Worked Direct-Product Example}

Take \(A=0\), a fixed compact \(X_6\), and the action~\eqref{eq:S10} with no
light scalar or vector zero modes.  Then
\[
 \sqrt{|g_{10}|}=\sqrt{-g_4}\sqrt{h_6},
\qquad
 R_{10}=R_4+R_6,
\]
and direct integration gives
\begin{align}
 S_{10}
 &=
 \frac{V_6}{2\kappa_{10}^2}
 \int_{Y_4}\sqrt{-g_4}\,
 \bigl(R_4+\overline R_6-2\Lambda_{10}\bigr)\,\dd^4x
 +S_{\rm matter}^{(4)}
\nonumber\\
 &=
 \frac{1}{2\kappa_4^2}
 \int_{Y_4}\sqrt{-g_4}\,
 (R_4-2\Lambda_4)\,\dd^4x
 +S_{\rm matter}^{(4)},
\end{align}
with
\[
 \kappa_4^2=\frac{\kappa_{10}^2}{V_6},
\qquad
 \Lambda_4=\Lambda_{10}-\frac12\overline R_6
 +\text{(matter/flux vacuum terms)}.
\]
This calculation is exact for the declared direct-product ansatz.  It is not
yet a proof that the ansatz solves the full ten-dimensional equations.
Equation~\eqref{eq:exacttrunc} supplies that additional test.

Now allow \(V_6=V_6(x)\).  The volume becomes a four-dimensional scalar, and a
Weyl transformation to Einstein frame generates its kinetic term and
couplings.  Pure GR is recovered only after this modulus is stabilized or
shown to decouple.  This simple comparison explains why ``integrate over the
fiber'' is not, by itself, a complete compactification proof.

\section{Claim Ledger and Research Frontier}

\begin{center}
\begin{tabular}{p{0.31\linewidth}p{0.17\linewidth}p{0.42\linewidth}}
\toprule
Claim & Status & Required evidence \\
\midrule
\(M_{10}\to Y_4\) with six-dimensional \(X_6\) &
Declared realization &
Global bundle, metric, connection, and selected compactification solution. \\
Einstein--Hilbert reduction &
Proved conditionally &
Action~\eqref{eq:S10}, zero-mode audit, and
Eq.~\eqref{eq:exacttrunc} or Theorem~\ref{thm:controlled}. \\
Einstein equation and Bianchi identity &
Proved conditionally &
Local diffeomorphism invariance, metric variation, and retained matter
equations. \\
Two TT polarizations on the exact finite branch &
Derived exactly &
The \(\mathbb Z_{64}\) support and intertwining certificate. \\
Physical \(G_N\) and Planck scale &
Open &
Selected \(\kappa_{10}\), normalized \(V_W\), pole residue, and unit map. \\
Full stress-energy response &
Open &
One normalized source operator coupling all retained matter to the TT and
constraint sectors. \\
Projection-only emergence of GR &
Open &
Selection of Lorentzian principal symbol, local action, and nonlinear
universal coupling from one upper source. \\
\bottomrule
\end{tabular}
\end{center}

The next decisive theorem is therefore not another dimensional-reduction
identity.  It is a same-source gravity theorem that constructs the Lorentzian
principal symbol, normalized spin-two kinetic operator, and conserved matter
response from the selected MTT geometry, while proving the zero-mode and
consistent-truncation conditions used here.

\section{Conclusion}

The corrected result is substantial but narrower than the previous title
suggested.  MTT currently has a valid conditional path from a
ten-dimensional metric realization to four-dimensional Einstein gravity and
an exact internal helicity-two support certificate.  The dimensional
reduction, effective Newton-coupling dictionary, Einstein variation, Bianchi
identity, and geometric-optics limit all follow once their stated inputs are
present.

What remains is the genuinely theory-defining step: selecting those inputs
from one MTT source and fixing their physical normalization.  Until that step
is complete, this paper establishes controlled compatibility with GR, not a
complete derivation of GR from projection alone.

% BEGIN MTT MANAGED COMPUTATIONAL EVIDENCE
\section*{Computational Evidence and Reproducibility}
The internal transverse-traceless certificate is used directly as support for the declared exact branch. Its own boundary is preserved: physical normalization remains open. The paper's four-dimensional Einstein reduction still depends on the stated coherent-reduction and action hypotheses.

The referenced rows are frozen to the curated results repository state identified below. Hashes are grouped in eight-character blocks for line breaking.
\begin{quote}\small
\textbf{Repository:} \url{https://github.com/PeterNero/mtt-results-repro}\\
\textbf{Commit:} \texttt{31247ebb 5c22f3fb b5443024 365433c6 ee0bff4a}\\
\textbf{Manifest:} \path{release/result_manifest.json}\\
\textbf{Manifest SHA-256:}\\
\texttt{fb399689 60b00584 631dbf53 1a708e18}\\
\texttt{ef928d6b 6d935119 c185d7f6 32b1e7cd}
\end{quote}

Tier labels are quoted verbatim from that manifest. A row used directly supports only the specific computational statement identified above; a corpus-state cross-check does not prove this paper's local theorems; and an open row is evidence of an unresolved obligation, never of closure.

\par\begingroup\dimen0=\pagegoal\advance\dimen0 by -\pagetotal\ifdim\dimen0<7\baselineskip\newpage\fi\endgroup
\paragraph{Rows used directly in this paper.}
\begin{itemize}
\setlength{\itemsep}{0.25em}
\item \path{A13/gr_tt_support} (\emph{derived exact}).\par Exact-branch internal TT support certificate; physical normalization remains open.
\end{itemize}

No imported row changes theorem ownership or promotes a neighboring claim: all local statements retain their stated hypotheses, domains, and limitations.
% END MTT MANAGED COMPUTATIONAL EVIDENCE

\begin{thebibliography}{99}

\bibitem{MTTFoundation}
P.~Nero,
\emph{Modal Triplet Theory: Foundations},
revised edition, 2026.

\bibitem{MTTAction}
P.~Nero,
\emph{Closure Geometry and a Regime-Local Ten-Dimensional Action Ansatz},
fourth edition, 2026.

\bibitem{MTTB1}
P.~Nero,
\emph{The Modal Triplet Theory Program B1: Loop-Transport Consistency and
the Conditional Gravity Realization},
second edition, 2026.

\bibitem{Lovelock}
D.~Lovelock,
``The Einstein tensor and its generalizations,''
\emph{Journal of Mathematical Physics} \textbf{12} (1971) 498--501,
\href{https://doi.org/10.1063/1.1665613}
{doi:10.1063/1.1665613}.

\bibitem{DuffKK}
M.~J.~Duff, B.~E.~W.~Nilsson, and C.~N.~Pope,
``Kaluza--Klein supergravity,''
\emph{Physics Reports} \textbf{130} (1986) 1--142,
\href{https://doi.org/10.1016/0370-1573(86)90163-8}
{doi:10.1016/0370-1573(86)90163-8}.

\bibitem{Petrini}
M.~Petrini,
``A systematic approach to consistent truncations of supergravity
theories,''
\emph{Universe} \textbf{7} (2021) 485,
\href{https://doi.org/10.3390/universe7120485}
{doi:10.3390/universe7120485}.

\bibitem{Donoghue}
J.~F.~Donoghue,
``General relativity as an effective field theory: The leading quantum
corrections,''
\emph{Physical Review D} \textbf{50} (1994) 3874--3888,
\href{https://doi.org/10.1103/PhysRevD.50.3874}
{doi:10.1103/PhysRevD.50.3874}.

\bibitem{Wald}
R.~M.~Wald,
\emph{General Relativity},
University of Chicago Press, 1984.

\bibitem{GibbonsHawking}
G.~W.~Gibbons and S.~W.~Hawking,
``Action integrals and partition functions in quantum gravity,''
\emph{Physical Review D} \textbf{15} (1977) 2752--2756,
\href{https://doi.org/10.1103/PhysRevD.15.2752}
{doi:10.1103/PhysRevD.15.2752}.

\end{thebibliography}

\end{document}
